\documentclass{nwafu-proposal}

\newcommand{\graduationyear}{2027}
\newcommand{\studentid}{}
\newcommand{\reporttitle}{请填写毕业论文（设计）题目}
\newcommand{\college}{}
\newcommand{\major}{}
\newcommand{\studentclass}{}
\newcommand{\studentname}{}
\newcommand{\supervisor}{}
\newcommand{\cosupervisor}{无}
\newcommand{\proposaldate}{2026年10月24日}

\begin{document}
\makeproposalcover
\makereporttitle

\section{选题依据}
\subsection{国内外研究现状}
请梳理本领域国内外研究进展，说明已有成果、主要问题和发展趋势。
引用采用著者—出版年制，例如\citet{huang2004}，或括号引用\citep{hull2002}。

\subsection{选题目的与意义}
请说明选题的理论意义、应用价值以及拟解决的问题。

\section{研究目标}
请明确研究目标，并给出可验证的评价指标。

\section{研究（设计）内容}
请分条阐述研究内容、重点问题和各项任务之间的关系。

\section{研究方案}
\subsection{研究方法}
请说明研究对象、数据来源、实验设计、分析方法和实施条件。
必要时可使用脚注补充说明。\footnote{此处为脚注示例，正式使用时请替换或删除。}

\subsection{技术路线}
\begin{figure}[htbp]
\centering
\fbox{\parbox[c][2.5cm][c]{.85\linewidth}{\centering 文献调研与问题分析 $\longrightarrow$ 方案设计\par 数据获取与分析 $\longrightarrow$ 结果验证与总结}}
\caption{研究技术路线示意图}
\label{fig:workflow}
\end{figure}
请结合图\ref{fig:workflow}说明研究步骤。
公式排版示例为
\begin{equation}
\bar{x}=\frac{1}{n}\sum_{i=1}^{n}x_i .
\end{equation}

\section{预期结果}
请说明预期形成的研究结论、设计成果、模型或其他可交付成果。

\section{存在困难}
请分析研究条件、实验材料、数据获取和技术实施方面的困难及解决办法。

\section{风险评估}
\subsection*{一、所用实验室类型}
\noindent
\riskcheck 化学类\quad \riskcheck 生物类\quad \riskcheck 机电类\quad \riskcheck 特种设备类\par
\noindent
\riskcheck 辐射安全类\quad \riskcheck 基础安全\quad \riskcheck 其他

\subsection*{二、实验项目使用的危险源及风险分析}
\noindent 危险源种类：\par
\noindent
\riskcheck 化学安全\quad \riskcheck 特种设备安全\quad \riskcheck 气瓶安全\quad \riskcheck 用电安全\par
\noindent
\riskcheck 生物安全\quad \riskcheck 辐射安全\quad \riskcheck 基础安全\quad \riskcheck 其他

\subsubsection*{危险源及风险分析}
请根据实验项目所使用的危险源列出具体清单，如管控类化学品名称、各种特殊设备名称等；分析实验过程中可能对人身安全、人体健康、实验室环境和周边环境带来的负面影响。

\subsubsection*{拟采取的防护和应急措施}
请对照《高校实验室安全检查项目表》、化学品安全技术说明书和相关管理制度，逐项说明防护措施及应急处理流程。

\subsubsection*{本科毕业论文（设计）执行人承诺}
本人对实验项目存在的风险进行全面了解，并认真分析，承诺在开展实验时做好个人防护，按照操作流程规范操作，并对实验过程的安全负责。

\noindent 执行人签字：\blankfield{3cm}{}\hfill 年\quad 月\quad 日\par

\subsubsection*{本科毕业论文（设计）指导老师承诺}
该生参加并通过了学校及学院组织的安全教育及培训，本人全面分析了实验项目中存在的风险，教育学生要注意实验安全，要求学生做好安全防护，按照操作流程进行规范操作，并负有安全管理责任。

\noindent 指导老师签字：\blankfield{3cm}{}\hfill 年\quad 月\quad 日


\section{进度安排}
\begin{table}[htbp]
\caption{研究进度安排}
\label{tab:schedule}
\centering
\fontsize{10.5pt}{18pt}\selectfont
\begin{tabularx}{\linewidth}{@{}p{3cm}XX@{}}
\toprule
时间 & 主要工作 & 阶段成果 \\
\midrule
起止日期 & 文献调研、明确研究问题 & 文献综述与开题报告 \\
起止日期 & 开展实验或设计、收集数据 & 实验记录与初步结果 \\
起止日期 & 数据分析、结果验证 & 分析结果与论文初稿 \\
起止日期 & 修改论文、准备答辩 & 论文定稿与答辩材料 \\
\bottomrule
\end{tabularx}
\end{table}

\section{指导教师意见（签名）}
\vspace*{2cm}
\begin{flushright}
指导教师签名：\blankfield{3cm}{}\par
日期：\blankfield{3cm}{}
\end{flushright}

\begin{thebibliography}{99}
\bibitem[黄骥等(2004)]{huang2004}
黄骥, 王建飞, 张红生. 2004. 植物戊糖磷酸途径及其两个关键酶的研究进展[J]. 植物学报, 21(2): 139-145.
\bibitem[Hull(2002)]{hull2002}
Hull R. 2002. Matthews' Plant Virology[M]. London: Academic Press.
\end{thebibliography}
\end{document}
